% Diagrama de cuerpo libre del nodo 3 del ejemplo A, con las fuerzas de las
% barras supuestas en tracción (capítulo 10)
\begin{tikzpicture}[scale=1.0,
    fuerza barra/.style={-{Stealth[length=7pt]}, line width=1.2pt,
      color=principal},
    texto/.style={font=\small}]
  \coordinate (n3) at (0,0);
  % Barras, recortadas
  \draw[barra, color=neutro, line width=1.2pt] (n3) -- ++(1.5,-2);
  \draw[barra, color=neutro, line width=1.2pt] (n3) -- ++(-1.5,-2);
  \node[texto, color=texto-suave] at (1.85,-2.15) {hacia el nodo 2};
  \node[texto, color=texto-suave] at (-1.85,-2.15) {hacia el nodo 1};
  % Fuerzas de las barras (tracción: tiran del nodo hacia la barra)
  \draw[fuerza barra] (n3) -- ++(0.9,-1.2) node[right, texto] {$N_2$};
  \draw[fuerza barra] (n3) -- ++(-0.9,-1.2) node[left, texto] {$N_3$};
  % Cargas
  \cargapuntual{(n3)}{(1.3,0)}{\SI{20}{\kilo\newton}}
  \cargapuntual{(n3)}{(0,-1.3)}{\SI{50}{\kilo\newton}}
  \node[nodo] at (n3) {};
  \node[etiqueta nodo, above right=3pt] at (n3) {3};
  % Direcciones
  \node[texto, anchor=west, align=left] at (2.6,0.2)
    {$N_2$: dirección $(0.6, -0.8)$\\ $N_3$: dirección $(-0.6, -0.8)$};
\end{tikzpicture}
